\documentclass[10pt,a4paper]{amsart}
\usepackage[margin=19mm]{geometry}
\usepackage{amsmath,amssymb,mathtools}
\usepackage[scr=rsfs,cal=euler]{mathalfa}
\usepackage{xfrac}
\usepackage[unicode,hidelinks,bookmarks=false]{hyperref}
\newcommand{\ie}{i.e.\ }
\newcommand{\cf}{cf.\ }
\newcommand{\resp}{respectively\ }
\newcommand{\Subsection}{\subsection}
\providecommand{\nobreakdash}{\nobreak\hbox{-}}
\setlength{\emergencystretch}{2em}
\multlinegap0pt
\usepackage{amssymb}
\newtheorem{corollary}{Corollary}
\newtheorem{lemma}{Lemma}
\newtheorem{theorem}{Theorem}
\title{Exact $6$-cut rigidity and small-order superconnectivity\\ for the $6$-regular case of Dirac's $k=4$ problem}
\author{Alper Ferudun}
\date{Source: arXiv:2606.18462v1 (CC BY 4.0)}
\begin{document}
\maketitle

\noindent\emph{Source and licence.} Exact $6$-cut rigidity and small-order superconnectivity\\ for the $6$-regular case of Dirac's $k=4$ problem. Version arXiv:2606.18462v1 (CC BY 4.0), distributed by arXiv under the Creative Commons Attribution 4.0 International licence, http://creativecommons.org/licenses/by/4.0/, as recorded in the arXiv licence metadata for this exact version and shown on the abstract page https://arxiv.org/abs/2606.18462v1. Full licence notice: \texttt{licenses/arxiv-2606.18462-CC-BY-4.0.txt}.\par\smallskip

\noindent\textit{Editorial scope.} Counted unit: the article's theorem thm:C (source0.tex lines 450--453). The article's complete original proof of it is delivered with it (source0.tex lines 455--476, one \texttt{proof} environment of 22 lines). No mathematics is written, repaired or re-derived: the statement and the proof are the article's own lines, in the article's own notation, and no retained line is substituted or re-flowed.  Retained uncounted as the source-local context the proof needs: lines 266--270; lines 402--404; lines 423--426.  Nothing was omitted from any retained span, so the package carries no omission record.  No retained line contains a citation, so the wrapper carries no bibliography and no bibliography entry.  No retained line contains a figure environment, an image-inclusion command or drawing code, so no image is delivered and the package declares no asset dependency.  Editorial formatting only, in addition: an AMS \texttt{amsart} preamble that restates the article's own declarations and macros used by the retained lines, namely the environments \texttt{corollary}, \texttt{lemma}, \texttt{theorem} on separate counters, together with the standard packages those lines need.  The retained environments print in this document as Corollary 1, Lemma 1, Theorem 1 in order of appearance, whereas the article prints the same items with the numbering of its own full document.  The complete native source of the article as distributed in the arXiv e-print for this exact version is bundled unchanged as \texttt{sources/203079/source0.tex}.
\begin{corollary}\label{cor:rowsum}
For every $3$-colouring $\alpha$ of a $6$-cut shore $G[A]$ in a target, the vector
$\bigl(\sum_{\alpha(v)=i} b(v)\bigr)_{i=1,2,3}$, where $b(v)$ is the number of cut
edges at $v$, is a permutation of $(6,0,0)$, $(3,3,0)$, $(4,1,1)$ or $(2,2,2)$.
\end{corollary}

\begin{lemma}[pair rigidity]\label{lem:pairrigid}
A concentrated shore of a target is never pair-rich.
\end{lemma}

\begin{lemma}[anti-diagonal exclusion]\label{lem:antidiag}
A concentrated shore of a target never satisfies
``$\alpha(p)\ne\alpha(\ell)$ for every proper $3$-colouring $\alpha$ of $H$''.
\end{lemma}

\begin{theorem}[C]\label{thm:C}
In a $6$-regular $(4,1)$-graph, no shore of a nontrivial $6$-edge-cut induces a
bipartite graph.
\end{theorem}

\begin{proof}
Let $H$ be bipartite with parts $X\cup Y$. First, the deficiency of $H$ must be
concentrated. Indeed, colouring $X$ with colour $1$ and an arbitrary subset
$T\subseteq Y$ with colour $2$, $Y\setminus T$ with colour $3$, is always proper;
by Corollary~\ref{cor:rowsum} every subset sum $t$ of the $Y$-side deficiencies
must keep $\bigl(b(X),t,b(Y)-t\bigr)$ among the allowed vectors. Combined with
the plain bipartition colouring (which forces
$(b(X),b(Y))\in\{(6,0),(3,3),(0,6)\}$) and the same argument on the $X$ side,
the only solutions with all $b(v)\le 5$ are: a single vertex of deficiency $3$
on each side, or two vertices of deficiency $3$ on one side.

If $p,\ell$ are nonadjacent, $H$ is pair-rich: for $p,\ell$ in the same part
$X$, colour $p$ and $X\setminus\{\ell\}$ with $1$, $\ell$ with $2$, $Y$ with $3$
(off-diagonal pairs; the bipartition colouring gives diagonal pairs); for
$p\in X$, $\ell\in Y$ nonadjacent, colour $p,\ell$ with $1$, $X\setminus\{p\}$
with $2$, $Y\setminus\{\ell\}$ with $3$ (diagonal pairs; the bipartition
colouring gives off-diagonal pairs). Colour permutations complete all nine
pairs, contradicting Lemma~\ref{lem:pairrigid}.

If $p,\ell$ are adjacent (necessarily in different parts), every colouring has
$\alpha(p)\ne\alpha(\ell)$, contradicting Lemma~\ref{lem:antidiag}.
\end{proof}

\end{document}
